\section{Overview}
\subsection{Relocatable object model}

The assembler does not assign final executable addresses. Each source file is assembled
into a relocatable object containing object-local bytes, symbol definitions, and
relocation records. Labels define object-local offsets. References to address-bearing
symbols are recorded for the linker when their final value depends on object placement.

The linker combines objects in the order supplied on the command line, assigns each
object a final base address, resolves non-mangled symbols, applies relocations, and
writes the final flat binary image.

\subsection{VIV object files}

The assembler does not emit a final executable image directly. Instead, it emits a
VIV object file, using the \texttt{.vo} format described in
%\cref{sec:toolchain-object-format}%
. A VIV object file contains object-local bytecode,
a symbol table, and a relocation table. The bytecode region is opaque and may contain
instructions, data, strings, and padding. Symbol and relocation records describe how
the linker should resolve symbolic references once final object placement is known.

\subsection{Linking}

The \texttt{viv32-ld} linker combines one or more VIV object files into a final flat
binary image. Input object files are laid out in the order supplied on the command
line. This ordering is significant: early objects may contain reset vectors,
interrupt vectors, or startup code that must appear before later program code. The
linker assigns each object a final base address, resolves non-mangled symbols, applies
relocation records, and writes the final executable byte stream.
